\documentclass[letterpaper]{article} % DO NOT CHANGE THIS
\usepackage[submission]{aaai2027}  % DO NOT CHANGE THIS
\usepackage[hyphens]{url}  % DO NOT CHANGE THIS
\usepackage{graphicx} % DO NOT CHANGE THIS
\urlstyle{rm} % DO NOT CHANGE THIS
\def\UrlFont{\rm}  % DO NOT CHANGE THIS
\usepackage{natbib}  % DO NOT CHANGE THIS AND DO NOT ADD ANY OPTIONS TO IT
\usepackage{caption} % DO NOT CHANGE THIS AND DO NOT ADD ANY OPTIONS TO IT
\frenchspacing  % DO NOT CHANGE THIS
\usepackage{amsmath}
\usepackage{amssymb}
\usepackage{booktabs}
\usepackage{enumitem}

\pdfinfo{
/TemplateVersion (2027.1)
}

\setcounter{secnumdepth}{2}

\title{Technical Appendix:\\ Motion Encoding for Human Perceptual Similarity}
\author{Anonymous Submission}
\affiliations{}

\begin{document}
\maketitle

\noindent This appendix expands the checkpoint-selection protocol, the
capacity-control experiment, and the linear/nonlinear decodability probes, and reports the Pearson
correlations that accompany the primary Spearman metrics. The capacity-control and probe analyses
provide the fuller evidence for the dissociation results summarized in the main paper; the Pearson
correlations serve as a consistency check on the Spearman-based conclusions.

% =====================================================================
\section{Checkpoint-Selection Protocol}
\label{app:checkpoint}

Every encoder is pretrained under its own self-supervised objective (reconstruction, masked motion
prediction, or masked prediction with the auxiliary pose head) to a fixed budget, with the full
trajectory checkpointed every 100 epochs. The nested cross-validation harness is invoked as
\begin{quote}\ttfamily\small
cv\_nested.py --k 5 --repeats 5 \textbackslash\\
\hspace*{1.5em}--n-triplet-seeds 5 --triplet-epochs 200
\end{quote}
producing, per outer fold, a three-way disjoint partition into TEST, SELECT, and TRAIN subsets. Raw
Spearman correlation is reported on the TEST fold ($n=25 = 5\text{ folds} \times 5\text{ repeats}$).

\paragraph{Which checkpoint is used.}
For each fold, the evaluated checkpoint is chosen from a grid of every 100th epoch within
$[100,\,\mathrm{PLATEAU}]$, where PLATEAU is the encoder's \emph{own} reconstruction-loss plateau
epoch. Selection therefore ranges only over checkpoints up to the point at which the model's
self-supervised reconstruction has converged. Within this window, the reported checkpoint is the one
that maximizes the mean raw Spearman correlation across the three actions on the SELECT fold, so that
\emph{perceptual} alignment---not reconstruction---determines which stored state is used. Two properties
make this choice leakage-clean:

\begin{itemize}[leftmargin=1.4em,itemsep=1pt]
\item PLATEAU is defined on the \emph{reconstruction} side --- the validation reconstruction loss
  (\texttt{val\_recon}, mean squared error in normalized space on a $5\%$ held-out split, seed $0$) ---
  never on perceptual performance. This keeps the SELECT fold from being used twice.
\item PLATEAU is \emph{model-intrinsic}, not the arbitrary end of a training trajectory (training
  length is a scheduling artifact). The SELECT fold only chooses \emph{within} the converged window;
  it does not set its boundary.
\end{itemize}

\paragraph{Operational definition of the plateau.}
The validation reconstruction curve is smoothed with a 3-point moving average. The asymptote is the
minimum over the final $20\%$ of the trajectory. PLATEAU is the first grid-100 epoch at which the
smoothed curve first falls to $\le 1.02\times$ the asymptote and remains $\le 1.03\times$ the asymptote
thereafter. The headline result is stable across the selection window (ceiling-invariance): it is
essentially unchanged whether checkpoints are drawn from $[0,\mathrm{PLATEAU}]$,
$[0,2\cdot\mathrm{PLATEAU}]$, or the full trajectory. The masked-prediction family can diverge if
over-trained; for these models PLATEAU is taken as the pre-divergence epoch under the same rule.

\paragraph{Folds, seeds, and hardware.}
Evaluation uses $5\text{ repeats}\times 5\text{ folds}$ nested cross-validation ($n=25$) with stratified
seeds $42 + 1000\cdot\text{repeat}$; folds shuffle only the canonical 57 Effort classes per action, and
perceptual fine-tuning runs for a fixed 200-epoch budget without early stopping on test sets.
Experiments ran on Rocky Linux 8.10 nodes with NVIDIA A100 (80\,GB) and H200 (141\,GB) GPUs, using
PyTorch 2.4.1, CUDA 12.1, and Python 3.10; TMR baselines used PyTorch 2.5.1 and PyTorch-Lightning 2.6.

% =====================================================================
\section{Capacity-Control Experiment}
\label{app:capacity}

To test whether the reconstruction backbone's perceptual advantage is merely a matter of model size, we
scaled the plain reconstruction transformer to match the MLD backbone's parameter budget (35.0M
parameters; hidden dimension $384$, $6$ attention heads, depth $7$; identical reconstruction objective
and pretraining corpus). Table~\ref{tab:capacity-supp} reports the outcome.

\begin{table}[htbp]
\centering
\small
\setlength{\tabcolsep}{5pt}
\caption{Capacity control: scaling the plain reconstruction transformer to the MLD budget. Larger
capacity lowers reconstruction loss to the best of any encoder we trained, yet does not improve
perceptual alignment.}
\label{tab:capacity-supp}
\begin{tabular}{lcc}
\toprule
Quantity & Small ($\approx$9M) & Large ($35.0$M) \\
\midrule
Val.\ reconstruction loss & $.00563$ & $\mathbf{.00305}$ \\
Pointing $\rho$ & $.213{\scriptstyle\pm}.034$ & $.221{\scriptstyle\pm}.037$ \\
Walking $\rho$ & $\mathbf{.476}$ & $.349$ \\
Picking $\rho$ & $\mathbf{.491}$ & $.301$ \\
\bottomrule
\end{tabular}
\end{table}

The larger model reconstructs best of any encoder we trained ($.00305$ versus $.00563$) and plateaus
faster, yet pointing alignment is statistically unchanged ($.221$ vs.\ $.213$) and walking and picking
decline as the model overfits the reconstruction target. A $3.9\times$ increase in parameters therefore
buys better reconstruction and zero perceptual gain, so the MLD-over-plain-transformer pointing gap
($\approx\!0.11$) is driven by architecture and objective rather than by parameter count.

% =====================================================================
\section{Decodability Probes}
\label{app:probe}

Section~8 of the main paper reports that, on pointing, perceptual alignment runs inversely to the
linear decodability of per-frame pose. This section gives the probe setup and the full nonlinear
control.

\paragraph{Linear (ridge) probe.}
For each frozen encoder we fit a ridge-regression probe from the mean-pooled embedding to per-frame
joint pose. Pose targets are reduced to $k=16$ PCA components and standardized per dimension; padded
frames are masked out; the ridge coefficient is set by generalized cross-validation. Decodability is
reported as $R^2$ on held-out frames under the same fold structure as the perceptual evaluation.

\paragraph{Nonlinear (MLP) control.}
To rule out a weak-probe artifact, we repeat the analysis with a multilayer-perceptron probe of the
same targets. The nonlinear probe recovers substantially more pose from every encoder (for MAMP+pose,
$R^2$ rises from $.509$ to $.645$; for the vanilla attention-pooling encoder, from $.727$ to $.829$)
while \emph{preserving the ordering} across encoders. Two conclusions follow: the inverse relationship
is not an artifact of linear probing, and per-frame pose is present in every encoder, including the
best-aligned one. The difference is that, in the encoders that align most closely with perception, pose
is arranged so as to be less \emph{linearly} accessible and therefore less visible to the $L_2$ metric
that both the learned and geometric distances use. A supporting control confirms that per-frame
velocity and acceleration are \emph{not} linearly decodable from any encoder ($R^2 < 0$ throughout), so
no encoder simply stores per-frame dynamics in linearly accessible form.

\paragraph{Action specificity.}
The inverse alignment-versus-decodability relationship is measurable only on pointing, the action on
which the encoders separate enough in both quantities. On walking and picking, linear pose decodability
is essentially flat across encoders and does not track perceptual alignment; we therefore make no
mechanistic claim for those actions.

% =====================================================================
\section{Pearson Correlation Analysis}
\label{app:pearson}
Table~\ref{tab:pearson-supp} reports Pearson correlation ($r$) for the main evaluation splits, as a
magnitude counterpart to the primary Spearman ($\rho$) results in the main paper. The Pearson metrics
preserve the same encoder ordering as those Spearman results, so the ordinal conclusions are unchanged
under a linear-correlation criterion.

\begin{table}[htbp]
\centering
\small
\setlength{\tabcolsep}{5pt}
\caption{Pearson correlation ($r$) across actions (Walking / Pointing / Picking, $n=25$).}
\label{tab:pearson-supp}
\begin{tabular}{lc}
\toprule
Encoder & Pearson $r$ (W / P / K) \\
\midrule
MAMP+pose & $.569$ / $.472$ / $.604$ \\
MLD AE & $.537$ / $.342$ / $.535$ \\
MAMP (Motion-only) & $.341$ / $.336$ / $.543$ \\
Vanilla AE & $.494$ / $.206$ / $.548$ \\
\bottomrule
\end{tabular}
\end{table}

\end{document}
